\section{Returning to the original linked field}

% BEGIN L69D1
\paragraph{L69D1. The quadratic character required by the rank-two run.}
The source predicate $\mathrm{Twist}$ concerns every imaginary
quadratic field $K$ unramified at the odd prime $\ell$, every
comparison $\iota$ and the specified complex conjugation $c$.
Take the Hecke character $Y$ of D66T1 with conductor dividing
$|\operatorname{disc}K|$, exponents $(-1,0)$ at the chosen
embedding and its conjugate, and one exponent $m_0>0$ prime to
$\ell$ for the images of all finite local units. Its construction
uses only the discriminant, the quadratic character, and
$\ell\nmid\operatorname{disc}K$. In particular these steps also
work when $\ell$ is inert in $K$. Apply the reciprocity formula
of D66T2, using all continuous embeddings of each completion at
$\ell$, to obtain $\kappa_K$. For an inert rational prime
$q\nmid\ell\operatorname{disc}K$, evaluation at its Frobenius
square gives $-\varepsilon(q)^{-1}$, where $\varepsilon(q)=q$.
By Chebotarev and continuity this gives
$\kappa_K(t^2)=-\varepsilon(t)^{-1}$ for every
$t\in G_{\mathbb Q}\setminus G_K$. Substitution of $t=sc$ gives
$\kappa_K(s)\kappa_K(csc^{-1})=\varepsilon(s)^{-1}$ for
$s\in G_K$. This argument requires no splitting of $\ell$.
Outside $\ell$, reciprocity identifies the images of local units
with inertia, giving trivial inertia off the discriminant and
$\kappa_K^{m_0}=1$ on all inertia there and at the discriminant
primes away from $\ell$. At the trivial valuation inertia is
trivial. These are the conjugation, square, ramification and
uniform-order fields of $\mathrm{IsTwist}$ on the same continuous
character, with the required algebraic Hecke witness. This proves
$\mathrm{Twist}$ at its original scope.
% END L69D1

% BEGIN L69D2
\paragraph{L69D2. Algebraic twists under products and restriction.}
In the inverse-arithmetic Artin convention of CIP1, the idelic norm
character corresponds to the full cyclotomic character. In fact a
uniformizer away from $\ell$ has idelic norm $q^{-1}$ and maps
to inverse arithmetic Frobenius, where cyclotomic also has value
$q^{-1}$. Its archimedean exponents are all one on the positive
component: the two complex embeddings contribute the squared
complex modulus. Its $\ell$-adic avatar is continuous by the
reciprocity formula of D66T2. Agreement on the Frobenius elements
gives the claimed equality by density. An open-kernel Galois
character factors through a finite abelian quotient by compactness.
Compose reciprocity with this quotient and apply $\iota$ to obtain
its finite-order Hecke character, of exponent zero on the positive
component. Powers of the norm times this character give the first
part of $\mathrm{AlgCyc}$. For the product and inverse clauses
use the product and inverse of the Hecke characters. Their exponent
vectors add and negate, respectively. The Artin map can be taken
to be the same in both witnesses since its prescribed Frobenius
values and the local and global reciprocity law fix the map.

For an extension $E/F$ in $\mathrm{AlgRes}$ pull the Hecke character
of $F$ back by the idele norm $N_{E/F}$. At an embedding $\tau$ of
$E$ the resulting exponent is $a(\tau|_F)$: the norm is the product
over the embeddings over each embedding of $F$. This calculation
also works on the positive archimedean components when some real
places become complex. Apply it to the local factors at $\ell$ in
the formula for the avatar. Functoriality of Artin reciprocity
for $N_{E/F}$ and restriction of Galois groups gives the restriction
of the original avatar. Equivalently, at the unramified places
outside $\ell$ the norm on uniformizers raises the character value
to the residue degree, exactly the Frobenius restriction rule.
Density and continuity identify the characters on $G_E$.
Thus the algebraic-character assertions used in C69C2 are supplied
for the original comparison and embeddings.
% END L69D2

% BEGIN L69D3
\paragraph{L69D3. Two actual seeds and two runs on the given link.}
Fix the data of $\mathrm{LinkedEndEAlt}$ including the original
$\mathrm{SeedClause}$, and assume the unitary modularity-ideal
theorem $\mathrm{ModUAlt}$. Name the quadratic field and character
triple supplied by the clause $K,\chi',\eta',\theta$.
The pointwise comparison of the sign characters gives
$G_K=\ker\eta$. The modulus $M$ is nonzero. Otherwise the
splitting condition for the trivial valuation would put the
specified complex conjugation in $\ker\eta$, whereas its image on
the frame has sign $-1$. Also $\Phi.d>3C_0+7$, so the sum defining
$\Phi.d$ has a place to choose. This gives the distinguished local
index for the two-run theorem.

The link says that relative inertia for $L/L^+$ is trivial at
the finite places and the places dividing $M$ split: inertia,
or respectively the indicated decomposition group, maps into
$G_L$ inside $G_{L^+}$. These are the source conditions
$\mathrm{UnrSplit}(L,M)$. Outside $M$, rational cyclotomic and
the character of modulus $N$ are both trivial on inertia since
$\ell N\mid M$. The given congruence therefore makes $\chi'$
congruent to one there. Its order is prime to $\ell$, so reduction
is injective on that group of roots of unity, and $\chi'$ is
actually one. Finally $2\ell\mid M$ by evenness of $N$.

Use L69D1 on this $K$ and L69D2 with C69C2 to obtain the two
specified twists, $a_3=\varepsilon_L^{-1}$ and
$a_2=\kappa_K|_{G_L}$. Fix the single extension of coefficients
from C69C1 for both. Take the seed of S69B5 at bound $2$, using
the preceding inertia tests and $\chi'\ne\eta'$, and the seed
of S69B3 at the same bound with this $a_2$. Their spacings can
differ. In rank three the weight list is $(m_3,0,-m_3)$ and in
rank two it is $(m_2,-m_2)$, precisely $m_i$ times the ordered
list of the frame in each rank. Use each seed in P69A5 for
\emph{every} realization of its rank. The character comparison
and congruence in P69A5 are those of this one seed clause.
The resulting assertion supplies the required specialization on
every ideal having the source point property.

C69C7 gives the realizations avoiding each requested nonzero
integer. The level and ordinary determinant ring are available
there on the residual data in question, and KX5 supplies the full
weight algebra at either rank. Apply A68M5 on this common extended
frame. Its two runs give $\mathrm{FinPt}$ there and both
$\mathrm{AutPtAlt}$ conclusions already quantified over the
original test DVRs. RTC24 descends the finite assertion along the
chosen coefficient inclusion with the prescribed residue map.
We have obtained on the original link $k$
\[
 \mathrm{FinPt}(c.O,\Phi,3),\qquad
 \mathrm{AutPtAlt}(k,3,\iota,\varepsilon_L^{-1}),\qquad
 \exists a_2\quad
 \mathrm{IsAlgChar}(\iota,a_2)\ \land\
 \mathrm{AutPtAlt}(k,2,\iota,a_2).
\]
This proves $\mathrm{ModUAlt}\Rightarrow\mathrm{LinkedEndEAlt}$
with the classical seed inputs specified in the seed chapter.
% END L69D3

% BEGIN L69D4
\paragraph{L69D4. The remaining arithmetic requirements on a marked host.}
For each original host take the field, sign, modulus $M=\ell N$,
and mark of AFE16 under its cohomological hypothesis. Its
$\mathrm{SeedClause}$ is already on that field. For the Gaussian
sign the ramified rational prime divides $2$, hence divides $M$.
The chosen valuation lies over $\ell$, and AFE16 supplies the
$\mathrm{Ok}$ and $\mathrm{Big}$ conditions. L69D3 supplies the
two assertions needed in M65A4. With the required polarized
solvable descent on the twisted rank-three representation, M65A4
gives the finite end and the automorphic end for this same mark,
hence the whole universal $\mathrm{ThmS6}_s$ implication. This is
a conditional reduction. Its arithmetic modularity premise can in
turn be supplied by HIA2--3 and N69E5 from the source Hida
hypotheses, the given cohomological inputs and the four remaining
patching inputs of N69E4. In that route KX3 provides local
reciprocity coordinates. In particular the two-rank construction
now uses the actual arithmetic seed points. General patching and
the polarized descent of the automorphy window retain their
individual mathematical requirements.
% END L69D4
